\documentclass[11pt]{article}

% Shared notes settings for Elements of AIML lectures (UPES).
% Keep this file free of \documentclass to allow reuse via \input.

\usepackage[utf8]{inputenc}
\usepackage[T1]{fontenc}
\usepackage[a4paper,margin=1in]{geometry}
\usepackage{hyperref}
\usepackage{xurl}
\usepackage{booktabs}
\usepackage{amsmath}
\usepackage{amssymb}
\usepackage{listings}
\usepackage{xcolor}
\usepackage{enumitem}
\usepackage{parskip}

% Code listings (general-purpose).
\lstset{
  basicstyle=\ttfamily\small,
  keywordstyle=\color{blue},
  commentstyle=\color{gray},
  stringstyle=\color{teal},
  showstringspaces=false,
  columns=fullflexible,
  frame=single,
  framerule=0.2pt,
  breaklines=true
}

\setlist[itemize]{topsep=0.3em,itemsep=0.2em}
\setlist[enumerate]{topsep=0.3em,itemsep=0.2em}

% Simple per-section source line.
\newcommand{\notesources}[1]{\par\noindent{\small\color{gray}\textit{Sources:} \sloppy #1\par}}

% Unit-wise source strings for this subject (used by slides/notes).
% Path is relative to the lecture `latex/` folder (the compilation CWD).
\IfFileExists{../../../_shared/aiml-sources.tex}{%
  % Source strings for Elements of AIML (used by slides + notes).

\newcommand{\AIMLRepoURL}{https://github.com/tali7c/Elements-of-AIML}

\newcommand{\AIMLLabSources}{%
Provost and Fawcett, \textit{Data Science for Business}.;%
McKinney, \textit{Python for Data Analysis}.;%
WEKA Documentation (Waikato Environment for Knowledge Analysis), accessed 2026-02-28.;%
Matplotlib and Seaborn documentation, accessed 2026-02-28.%
}

\newcommand{\AIMLUnitISources}{%
Rich and Knight, \textit{Artificial Intelligence}, McGraw Hill, 2017.;%
Russell and Norvig, \textit{Artificial Intelligence: A Modern Approach} (4e), Ch. 1--2 (Introduction, Intelligent Agents).;%
Poole and Mackworth, \textit{Artificial Intelligence: Foundations of Computational Agents} (2e), selected sections.%
}

\newcommand{\AIMLUnitIISources}{%
Russell and Norvig, \textit{Artificial Intelligence: A Modern Approach} (4e), Ch. 7--9 (Logical Agents, First-Order Logic, Inference).;%
Huth and Ryan, \textit{Logic in Computer Science} (2e), selected sections on propositional and first-order logic.;%
Genesereth and Nilsson, \textit{Logical Foundations of Artificial Intelligence}, selected chapters.%
}

\newcommand{\AIMLUnitIIISources}{%
Mueller and Massaron, \textit{Machine Learning for Dummies}, For Dummies, 2016.;%
Mitchell, \textit{Machine Learning}, selected chapters on learning basics and evaluation.;%
Scikit-learn documentation, accessed 2026-02-28.%
}

\newcommand{\AIMLUnitIVSources}{%
Mueller and Massaron, \textit{Machine Learning for Dummies}, For Dummies, 2016.;%
James, Witten, Hastie, and Tibshirani, \textit{An Introduction to Statistical Learning} (2e), selected chapters.;%
Scikit-learn documentation, accessed 2026-02-28.%
}

\newcommand{\AIMLUnitVSources}{%
NIST, \textit{AI Risk Management Framework (AI RMF 1.0)}, 2023.;%
Barocas, Hardt, and Narayanan, \textit{Fairness and Machine Learning} (online book), selected sections.;%
Knaflic, \textit{Storytelling with Data}, selected chapters on communicating results.%
}
%
}{%
  % Faculty copies live under `private/` so their compile CWD is different.
  \IfFileExists{../../../../public/_shared/aiml-sources.tex}{%
    % Source strings for Elements of AIML (used by slides + notes).

\newcommand{\AIMLRepoURL}{https://github.com/tali7c/Elements-of-AIML}

\newcommand{\AIMLLabSources}{%
Provost and Fawcett, \textit{Data Science for Business}.;%
McKinney, \textit{Python for Data Analysis}.;%
WEKA Documentation (Waikato Environment for Knowledge Analysis), accessed 2026-02-28.;%
Matplotlib and Seaborn documentation, accessed 2026-02-28.%
}

\newcommand{\AIMLUnitISources}{%
Rich and Knight, \textit{Artificial Intelligence}, McGraw Hill, 2017.;%
Russell and Norvig, \textit{Artificial Intelligence: A Modern Approach} (4e), Ch. 1--2 (Introduction, Intelligent Agents).;%
Poole and Mackworth, \textit{Artificial Intelligence: Foundations of Computational Agents} (2e), selected sections.%
}

\newcommand{\AIMLUnitIISources}{%
Russell and Norvig, \textit{Artificial Intelligence: A Modern Approach} (4e), Ch. 7--9 (Logical Agents, First-Order Logic, Inference).;%
Huth and Ryan, \textit{Logic in Computer Science} (2e), selected sections on propositional and first-order logic.;%
Genesereth and Nilsson, \textit{Logical Foundations of Artificial Intelligence}, selected chapters.%
}

\newcommand{\AIMLUnitIIISources}{%
Mueller and Massaron, \textit{Machine Learning for Dummies}, For Dummies, 2016.;%
Mitchell, \textit{Machine Learning}, selected chapters on learning basics and evaluation.;%
Scikit-learn documentation, accessed 2026-02-28.%
}

\newcommand{\AIMLUnitIVSources}{%
Mueller and Massaron, \textit{Machine Learning for Dummies}, For Dummies, 2016.;%
James, Witten, Hastie, and Tibshirani, \textit{An Introduction to Statistical Learning} (2e), selected chapters.;%
Scikit-learn documentation, accessed 2026-02-28.%
}

\newcommand{\AIMLUnitVSources}{%
NIST, \textit{AI Risk Management Framework (AI RMF 1.0)}, 2023.;%
Barocas, Hardt, and Narayanan, \textit{Fairness and Machine Learning} (online book), selected sections.;%
Knaflic, \textit{Storytelling with Data}, selected chapters on communicating results.%
}
%
  }{%
    % Source strings for Elements of AIML (used by slides + notes).

\newcommand{\AIMLRepoURL}{https://github.com/tali7c/Elements-of-AIML}

\newcommand{\AIMLLabSources}{%
Provost and Fawcett, \textit{Data Science for Business}.;%
McKinney, \textit{Python for Data Analysis}.;%
WEKA Documentation (Waikato Environment for Knowledge Analysis), accessed 2026-02-28.;%
Matplotlib and Seaborn documentation, accessed 2026-02-28.%
}

\newcommand{\AIMLUnitISources}{%
Rich and Knight, \textit{Artificial Intelligence}, McGraw Hill, 2017.;%
Russell and Norvig, \textit{Artificial Intelligence: A Modern Approach} (4e), Ch. 1--2 (Introduction, Intelligent Agents).;%
Poole and Mackworth, \textit{Artificial Intelligence: Foundations of Computational Agents} (2e), selected sections.%
}

\newcommand{\AIMLUnitIISources}{%
Russell and Norvig, \textit{Artificial Intelligence: A Modern Approach} (4e), Ch. 7--9 (Logical Agents, First-Order Logic, Inference).;%
Huth and Ryan, \textit{Logic in Computer Science} (2e), selected sections on propositional and first-order logic.;%
Genesereth and Nilsson, \textit{Logical Foundations of Artificial Intelligence}, selected chapters.%
}

\newcommand{\AIMLUnitIIISources}{%
Mueller and Massaron, \textit{Machine Learning for Dummies}, For Dummies, 2016.;%
Mitchell, \textit{Machine Learning}, selected chapters on learning basics and evaluation.;%
Scikit-learn documentation, accessed 2026-02-28.%
}

\newcommand{\AIMLUnitIVSources}{%
Mueller and Massaron, \textit{Machine Learning for Dummies}, For Dummies, 2016.;%
James, Witten, Hastie, and Tibshirani, \textit{An Introduction to Statistical Learning} (2e), selected chapters.;%
Scikit-learn documentation, accessed 2026-02-28.%
}

\newcommand{\AIMLUnitVSources}{%
NIST, \textit{AI Risk Management Framework (AI RMF 1.0)}, 2023.;%
Barocas, Hardt, and Narayanan, \textit{Fairness and Machine Learning} (online book), selected sections.;%
Knaflic, \textit{Storytelling with Data}, selected chapters on communicating results.%
}
%
  }%
}


\title{Elements of AIML Lab\\Experiment 02 (After-submission Reference): WEKA for Classification}
\author{Tofik Ali}
\date{\today}

\begin{document}
\maketitle
\tableofcontents

\section{How to use this reference}
This document is meant to be shared \textbf{after you submit} Experiment~02.
Use it to verify that you:
\begin{itemize}[leftmargin=*]
  \item selected a \textbf{nominal} class attribute (label),
  \item evaluated models fairly (same split/CV),
  \item interpreted the confusion matrix and precision/recall/F1 correctly.
\end{itemize}

\section{Quick recap: what classification means}
Classification is \textbf{supervised learning} where the target is a \textbf{label} (nominal), e.g., spam/ham, pass/fail, disease/no-disease.

\subsection{Confusion matrix (binary)}
\begin{center}
\begin{tabular}{c|cc}
  & Pred + & Pred - \\\hline
  True + & TP & FN \\
  True - & FP & TN \\
\end{tabular}
\end{center}
Interpretation in plain words:
\begin{itemize}[leftmargin=*]
  \item TP: correct ``positive'' predictions
  \item FP: false alarms
  \item FN: missed positives
  \item TN: correct ``negative'' predictions
\end{itemize}

\subsection{Metrics (why more than accuracy)}
\begin{itemize}[leftmargin=*]
  \item Accuracy $= \frac{TP+TN}{TP+TN+FP+FN}$ (can be misleading if classes are imbalanced).
  \item Precision $= \frac{TP}{TP+FP}$ (quality of positive predictions).
  \item Recall $= \frac{TP}{TP+FN}$ (coverage of true positives).
  \item F1 balances precision and recall.
\end{itemize}

\section{Worked example in WEKA (end-to-end)}
\subsection{Dataset choice (example)}
WEKA often provides sample classification datasets under a \texttt{data/} folder (e.g., \texttt{iris.arff}, \texttt{weather.nominal.arff}).
Your screenshots/metrics may differ depending on dataset choice.

\subsection{Step 1: load data and set class attribute}
\begin{enumerate}[leftmargin=*]
  \item WEKA $\rightarrow$ \textbf{Explorer} $\rightarrow$ \textbf{Preprocess} $\rightarrow$ \textbf{Open file...}
  \item Set \textbf{Class} to the label column.
\end{enumerate}
\textbf{Cross-check:} the class attribute must be \textbf{nominal}. If numeric, you are doing regression.

\subsection{Step 2: (optional) Replace missing values and remove IDs}
\begin{itemize}[leftmargin=*]
  \item Missing values: \texttt{Filter} $\rightarrow$ \texttt{unsupervised/attribute/ReplaceMissingValues}.
  \item Remove ID-like columns (roll number, customer ID) that leak identity rather than signal.
\end{itemize}

\subsection{Step 3: train and evaluate two classifiers}
\paragraph{Model A: Decision tree (J48).}
\begin{enumerate}[leftmargin=*]
  \item \textbf{Classify} tab $\rightarrow$ \textbf{Choose} $\rightarrow$ \texttt{trees/J48}.
  \item \textbf{Test options} $\rightarrow$ \textbf{Cross-validation} (10 folds) \emph{or} \textbf{Percentage split}.
  \item Click \textbf{Start}.
\end{enumerate}

\paragraph{Model B: Naive Bayes (a fast baseline).}
\begin{enumerate}[leftmargin=*]
  \item \textbf{Choose} $\rightarrow$ \texttt{bayes/NaiveBayes}.
  \item Keep the \textbf{same} test option as Model~A.
  \item Click \textbf{Start}.
\end{enumerate}

\subsection{Step 4: interpret WEKA output}
In the output, you typically see:
\begin{itemize}[leftmargin=*]
  \item overall accuracy,
  \item confusion matrix,
  \item per-class precision/recall/F1,
  \item weighted averages (important when classes are imbalanced).
\end{itemize}
\textbf{Practical reading tip:} If FN is costly (e.g., missing a disease), focus on \textbf{recall}. If FP is costly (e.g., blocking a genuine payment), focus on \textbf{precision}.

\section{What a strong submission looks like (cross-check)}
\subsection{Screenshots}
\begin{itemize}[leftmargin=*]
  \item Preprocess: correct class selected.
  \item Classify: model name + evaluation method.
  \item Output: confusion matrix + precision/recall/F1 for both models.
\end{itemize}

\subsection{Report structure (1--2 pages)}
\begin{enumerate}[leftmargin=*]
  \item Dataset and class labels (what does ``positive'' mean?).
  \item Preprocessing steps (if any).
  \item Two models, same evaluation method.
  \item Results table (accuracy, precision, recall, F1).
  \item Observations (trade-offs: which model increases recall but decreases precision, etc.).
\end{enumerate}

\section{Troubleshooting (common issues and fixes)}
\begin{itemize}[leftmargin=*]
  \item \textbf{Wrong class attribute:} selecting an ID column creates meaningless results.
  \item \textbf{Different evaluation settings across models:} comparisons become unfair.
  \item \textbf{Only reporting accuracy:} always include confusion matrix + precision/recall/F1.
  \item \textbf{Imbalance ignored:} if one class is rare, weighted averages and minority-class recall matter most.
\end{itemize}

\notesources{\AIMLLabSources}
\end{document}

